% ============================================================
\chapter{Experiment Tracking --- MLflow, Weights \& Biases}
\label{ch:experiment-tracking}
\index{experiment tracking}
% ============================================================

\section{Why Track Experiments?}
\label{sec:why-tracking}

During an ML project, a data scientist may run hundreds of experiments
varying hyperparameters, features, data splits, and model architectures.
Without systematic tracking, critical questions become unanswerable:

\begin{itemize}
  \item Which hyperparameters produced the best F1 score?
  \item What data version was used for model v2.3?
  \item Why did accuracy drop between Tuesday and Wednesday?
  \item Can we reproduce the result from last month?
\end{itemize}

\begin{definition}[Experiment tracking]
\textbf{Experiment tracking}\index{experiment tracking} is the practice of
systematically recording the parameters, metrics, artefacts, code versions,
and environment details of every ML experiment, enabling comparison,
reproducibility, and auditability.
\end{definition}

\begin{warning}[title=\textbf{The spreadsheet trap}]
Tracking experiments in spreadsheets or notebooks is fragile:
\begin{itemize}
  \item Manual entry is error-prone and incomplete.
  \item No link between logged values and actual code/data.
  \item No automatic artefact storage.
  \item Impossible to search, filter, or visualise at scale.
\end{itemize}
\end{warning}

\section{Experiment Tracking Architecture}
\label{sec:tracking-architecture}

\begin{center}
\begin{tikzpicture}[
  box/.style={rectangle, draw, rounded corners, minimum width=2.8cm,
    minimum height=1.2cm, align=center, font=\small},
  storage/.style={cylinder, draw, shape border rotate=90,
    minimum width=2cm, minimum height=1cm, align=center,
    font=\scriptsize, aspect=0.3},
  arr/.style={-{Stealth[length=2.5mm]}, thick, steelblue}
]
  % Training script
  \node[box, fill=blue!10] (script) {\textbf{Training Script}\\
    \scriptsize\texttt{train.py}};

  % Tracking server
  \node[box, fill=green!10, right=3cm of script] (server)
    {\textbf{Tracking Server}\\
    \scriptsize MLflow / W\&B};

  % Backend store
  \node[storage, fill=orange!10, above right=1.5cm and 1cm of server]
    (backend) {\textbf{Backend}\\Params, metrics};

  % Artefact store
  \node[storage, fill=red!10, below right=1.5cm and 1cm of server]
    (artefact) {\textbf{Artefact Store}\\Models, plots\\(S3, GCS)};

  % UI
  \node[box, fill=yellow!10, right=3cm of server] (ui)
    {\textbf{Web UI}\\Compare, visualise};

  % Arrows
  \draw[arr] (script) -- (server)
    node[midway, above, font=\scriptsize] {log params, metrics};
  \draw[arr] (server) -- (backend);
  \draw[arr] (server) -- (artefact);
  \draw[arr] (server) -- (ui);
\end{tikzpicture}
\end{center}

\section{MLflow}
\label{sec:mlflow}
\index{MLflow}

\textbf{MLflow}\index{MLflow} is an open-source platform for managing the
ML lifecycle. Its \emph{Tracking} component logs parameters, metrics, and
artefacts. Its \emph{Model Registry} manages model versions and deployment
stages.

\subsection{MLflow Tracking API}

\begin{codeblock}[title=\textbf{Complete MLflow tracking example}]
\begin{minted}{python}
import mlflow
import mlflow.sklearn
from sklearn.ensemble import RandomForestClassifier
from sklearn.datasets import load_iris
from sklearn.model_selection import train_test_split
from sklearn.metrics import accuracy_score, f1_score

# Set the tracking URI (local or remote server)
mlflow.set_tracking_uri("http://localhost:5000")
mlflow.set_experiment("iris-classification")

# Load data
X, y = load_iris(return_X_y=True)
X_train, X_test, y_train, y_test = train_test_split(
    X, y, test_size=0.2, random_state=42
)

# Experiment parameters
params = {
    "n_estimators": 200,
    "max_depth": 10,
    "min_samples_split": 5,
    "random_state": 42,
}

with mlflow.start_run(run_name="rf-baseline"):
    # Log parameters
    mlflow.log_params(params)

    # Train
    model = RandomForestClassifier(**params)
    model.fit(X_train, y_train)

    # Evaluate
    y_pred = model.predict(X_test)
    acc = accuracy_score(y_test, y_pred)
    f1 = f1_score(y_test, y_pred, average="weighted")

    # Log metrics
    mlflow.log_metric("accuracy", acc)
    mlflow.log_metric("f1_weighted", f1)

    # Log the model as an artefact
    mlflow.sklearn.log_model(model, "random-forest-model")

    # Log additional artefacts (plots, configs, etc.)
    mlflow.log_artifact("configs/config.yaml")

    print(f"Run ID: {mlflow.active_run().info.run_id}")
    print(f"Accuracy: {acc:.4f}, F1: {f1:.4f}")
\end{minted}
\end{codeblock}

\begin{codeblock}[title=\textbf{Launching the MLflow UI}]
\begin{minted}{bash}
# Start the tracking server
mlflow server --host 0.0.0.0 --port 5000 \
    --backend-store-uri sqlite:///mlflow.db \
    --default-artifact-root ./mlartifacts

# Open in browser: http://localhost:5000
\end{minted}
\end{codeblock}

\subsection{MLflow Model Registry}
\index{MLflow!model registry}

The Model Registry provides a centralised store for managing model versions
and their lifecycle stages (\emph{Staging}, \emph{Production},
\emph{Archived}).

\begin{codeblock}[title=\textbf{MLflow Model Registry workflow}]
\begin{minted}{python}
import mlflow
from mlflow.tracking import MlflowClient

client = MlflowClient()

# Register a model from a run
model_uri = f"runs:/{run_id}/random-forest-model"
result = mlflow.register_model(model_uri, "IrisClassifier")
print(f"Version: {result.version}")

# Transition to staging
client.transition_model_version_stage(
    name="IrisClassifier",
    version=result.version,
    stage="Staging",
)

# Promote to production
client.transition_model_version_stage(
    name="IrisClassifier",
    version=result.version,
    stage="Production",
)

# Load a production model for inference
model = mlflow.pyfunc.load_model(
    "models:/IrisClassifier/Production"
)
predictions = model.predict(X_test)
\end{minted}
\end{codeblock}

\begin{bestpractice}[title=\textbf{MLflow best practices}]
\begin{itemize}
  \item Use a \textbf{remote tracking server} (not local \file{mlruns/})
    for team collaboration.
  \item Use \textbf{experiment names} that reflect the project and task.
  \item Log \textbf{all} hyperparameters, not just the ones you are tuning.
  \item Use \cli{mlflow.autolog()} for supported frameworks (scikit-learn,
    PyTorch, TensorFlow) to automatically log parameters and metrics.
  \item Store artefacts (models, plots) in a shared object store (S3, GCS).
\end{itemize}
\end{bestpractice}

\section{Weights \& Biases (W\&B)}
\label{sec:wandb}
\index{Weights \& Biases}

\textbf{Weights \& Biases}\index{W\&B} (W\&B) is a cloud-hosted (or
self-hosted) experiment tracking platform with rich visualisation,
hyperparameter sweeps, and team collaboration features.

\begin{codeblock}[title=\textbf{Complete W\&B tracking example}]
\begin{minted}{python}
import wandb
import torch
import torch.nn as nn
from torch.utils.data import DataLoader, TensorDataset

# Initialise a W&B run
wandb.init(
    project="mnist-classification",
    name="cnn-baseline",
    config={
        "learning_rate": 1e-3,
        "batch_size": 64,
        "epochs": 20,
        "architecture": "SimpleCNN",
        "optimizer": "Adam",
    },
)
config = wandb.config

# Define model
class SimpleCNN(nn.Module):
    def __init__(self):
        super().__init__()
        self.conv1 = nn.Conv2d(1, 32, 3, padding=1)
        self.conv2 = nn.Conv2d(32, 64, 3, padding=1)
        self.pool = nn.MaxPool2d(2)
        self.fc = nn.Linear(64 * 7 * 7, 10)
        self.relu = nn.ReLU()

    def forward(self, x):
        x = self.pool(self.relu(self.conv1(x)))
        x = self.pool(self.relu(self.conv2(x)))
        x = x.view(x.size(0), -1)
        return self.fc(x)

model = SimpleCNN()
optimizer = torch.optim.Adam(
    model.parameters(), lr=config.learning_rate
)
criterion = nn.CrossEntropyLoss()

# Watch model (log gradients and parameters)
wandb.watch(model, log="all", log_freq=100)

# Training loop
for epoch in range(config.epochs):
    model.train()
    total_loss = 0
    correct = 0
    total = 0

    for X_batch, y_batch in train_loader:
        optimizer.zero_grad()
        outputs = model(X_batch)
        loss = criterion(outputs, y_batch)
        loss.backward()
        optimizer.step()

        total_loss += loss.item()
        _, predicted = outputs.max(1)
        correct += predicted.eq(y_batch).sum().item()
        total += y_batch.size(0)

    # Log metrics per epoch
    wandb.log({
        "epoch": epoch,
        "train/loss": total_loss / len(train_loader),
        "train/accuracy": correct / total,
    })

# Log final model as artefact
artifact = wandb.Artifact("cnn-model", type="model")
torch.save(model.state_dict(), "model.pt")
artifact.add_file("model.pt")
wandb.log_artifact(artifact)

wandb.finish()
\end{minted}
\end{codeblock}

\begin{remark}
W\&B automatically captures Git commit hashes, system information (GPU,
OS, Python version), and console output. The web dashboard allows
real-time comparison of runs with interactive charts.
\end{remark}

\section{Hydra for Configuration Management}
\label{sec:hydra}
\index{Hydra}

\textbf{Hydra}\index{Hydra} (by Meta) is a configuration framework that
simplifies managing complex configs with composition, overrides, and
multi-run sweeps.

\begin{codeblock}[title=\textbf{Hydra configuration example}]
\begin{minted}{yaml}
# configs/config.yaml
defaults:
  - model: random_forest
  - data: iris
  - _self_

training:
  seed: 42
  test_size: 0.2

# configs/model/random_forest.yaml
n_estimators: 200
max_depth: 10

# configs/model/gradient_boosting.yaml
n_estimators: 300
learning_rate: 0.1
max_depth: 5

# configs/data/iris.yaml
path: data/iris.csv
target_col: species
\end{minted}
\end{codeblock}

\begin{codeblock}[title=\textbf{Using Hydra in Python}]
\begin{minted}{python}
import hydra
from omegaconf import DictConfig
import mlflow

@hydra.main(version_base=None, config_path="configs",
            config_name="config")
def train(cfg: DictConfig) -> float:
    mlflow.set_experiment("hydra-experiments")

    with mlflow.start_run():
        mlflow.log_params(dict(cfg.model))
        # ... training code ...
        mlflow.log_metric("accuracy", accuracy)

    return accuracy

if __name__ == "__main__":
    train()
\end{minted}
\end{codeblock}

\begin{codeblock}[title=\textbf{Hydra command-line overrides and sweeps}]
\begin{minted}{bash}
# Override a parameter
python train.py model.n_estimators=500

# Switch model configuration
python train.py model=gradient_boosting

# Multi-run sweep
python train.py --multirun \
    model.n_estimators=100,200,500 \
    model.max_depth=5,10,20
\end{minted}
\end{codeblock}

\section{Tool Comparison}
\label{sec:tracking-comparison}

\begin{center}
\begin{tabular}{@{}lcccc@{}}
\toprule
\textbf{Criterion} & \textbf{MLflow} & \textbf{W\&B} & \textbf{Neptune} & \textbf{ClearML} \\
\midrule
Open source     & Yes   & Partial & Partial & Yes \\
Self-hosted     & Yes   & Yes     & Yes     & Yes \\
Cloud hosted    & No    & Yes     & Yes     & Yes \\
Free tier       & N/A   & Yes     & Yes     & Yes \\
Model registry  & Yes   & Yes     & Yes     & Yes \\
HPO sweeps      & No    & Yes     & Yes     & Yes \\
Real-time logs  & No    & Yes     & Yes     & Yes \\
Team collab     & Basic & Excellent & Good  & Good \\
Framework support & Broad & Broad  & Broad  & Broad \\
Learning curve  & Low   & Low     & Low     & Medium \\
\bottomrule
\end{tabular}
\end{center}

\begin{remark}
\textbf{MLflow} is the de facto standard for self-hosted, open-source
experiment tracking. \textbf{W\&B} excels in visualisation and team
collaboration for cloud-hosted workflows. Many teams use both: MLflow for
the model registry and W\&B for experiment visualisation.
\end{remark}

\section{Best Practices}

\begin{bestpractice}[title=\textbf{Experiment tracking golden rules}]
\begin{enumerate}
  \item \textbf{Log everything}: parameters, metrics, data version, code
    version, environment, random seeds.
  \item \textbf{Use meaningful run names}: \texttt{rf-200trees-depth10},
    not \texttt{run-42}.
  \item \textbf{Organise experiments}: one experiment per task or project.
  \item \textbf{Store artefacts}: save models, plots, and configs alongside
    metrics.
  \item \textbf{Automate logging}: use \cli{mlflow.autolog()} or W\&B
    integrations.
  \item \textbf{Compare systematically}: use the UI to compare runs, not
    memory or spreadsheets.
\end{enumerate}
\end{bestpractice}

\begin{antipattern}[title=\textbf{Common anti-patterns}]
\begin{itemize}
  \item Logging only the final metric, not per-epoch curves
  \item Forgetting to log the data version used for training
  \item Using \texttt{print()} instead of a tracking tool
  \item Running experiments without setting random seeds
  \item Not saving the model artefact alongside the metrics
\end{itemize}
\end{antipattern}

\section{Mini-project: Tracked Experiment Campaign}
\label{sec:mini-project-ch4}

\begin{miniproject}[title=\textbf{Mini-project 4: Full experiment tracking}]
Building on the DVC-versioned project from Chapter~\ref{ch:version-control}:
\begin{enumerate}
  \item Set up an MLflow tracking server (local or Docker).
  \item Instrument your training script to log parameters, metrics, and
    the trained model.
  \item Run at least 5 experiments varying hyperparameters.
  \item Use the MLflow UI to identify the best run.
  \item Register the best model in the MLflow Model Registry and promote
    it to \emph{Production}.
  \item Add W\&B tracking alongside MLflow and compare the dashboards.
  \item Use Hydra to manage experiment configurations and run a sweep
    over 3 hyperparameters.
\end{enumerate}
\end{miniproject}

\section{Exercises}
\label{sec:exercises-ch4}

\begin{exercise}[$\bigstar$ --- MLflow basics]
Write a training script that logs at least 5 parameters and 3 metrics to
MLflow. Launch the MLflow UI and take a screenshot of the comparison view
for three different runs. Explain what each column represents.
\end{exercise}

\begin{exercise}[$\bigstar\bigstar$ --- W\&B integration]
Rewrite the training script from the previous exercise using W\&B instead
of MLflow. Log training curves (loss and accuracy per epoch), a confusion
matrix, and a model artefact. Use \cli{wandb.sweep} to perform a
hyperparameter search over learning rate, batch size, and number of layers.
Compare the W\&B and MLflow experiences.
\end{exercise}

\begin{exercise}[$\bigstar\bigstar\bigstar$ --- End-to-end tracked pipeline]
Build a complete tracked pipeline for image classification (e.g., CIFAR-10):
\begin{enumerate}
  \item Use Hydra for configuration management with at least 3 config
    groups (model, data, training)
  \item Log all experiments to both MLflow and W\&B
  \item Implement early stopping based on validation loss
  \item Log per-epoch metrics, learning rate schedules, and sample
    predictions
  \item Use the MLflow Model Registry to manage model versions
  \item Write a comparison report using data exported from the tracking
    tools
\end{enumerate}
\end{exercise}

\section{Cheatsheet}

\begin{cheatsheet}[title=\textbf{Chapter 4 Summary}]
\begin{tabular}{@{}p{4cm}p{8.5cm}@{}}
\toprule
\textbf{Tool / Concept} & \textbf{Key Commands / API} \\
\midrule
\textbf{MLflow tracking} & \cli{mlflow.start\_run()},
  \cli{mlflow.log\_params()}, \cli{mlflow.log\_metric()} \\
\textbf{MLflow registry} & \cli{mlflow.register\_model()},
  \cli{client.transition\_model\_version\_stage()} \\
\textbf{MLflow server} & \cli{mlflow server --host 0.0.0.0 --port 5000} \\
\textbf{W\&B} & \cli{wandb.init()}, \cli{wandb.log()},
  \cli{wandb.log\_artifact()} \\
\textbf{Hydra} & \cli{@hydra.main()}, \cli{--multirun},
  \cli{model=gradient\_boosting} \\
\textbf{Best practice} & Log everything; meaningful names; store artefacts \\
\bottomrule
\end{tabular}
\end{cheatsheet}
